\section*{7 Binary Multiselection and Proof of Main Theorem (author source p.11:16)}
Next, we use the multiselection circuit over large alphabets to obtain a binary multiselection circuit.
\begin{originalresult}{Lemma 27}
Let $q=q(n)$ and $s=s(n)$ be integer functions of $n$, let $\Sigma=\bits^s$, and let $C_n$ be a (log-space uniform) Boolean circuit for $\Sel_\Sigma^{n/s\to q}$ with size $S$ and depth $D$.
Then there is a (log-space uniform) Boolean circuit $C_n'$ for $\Sel_{\{0,1\}}^{n\to q}$ with size $S+O(q\cdot s)$ and depth $D+O(\log s)$.
\end{originalresult}
\begin{proof}
On data input $x\in\bits^n$ and selector input $(i_1,\ldots,i_q)\in[n]^q$, $C_n'$ performs the following steps:
\begin{enumerate}
\item View $x\in\bits^n$ as $X\in\Sigma^{n/s}$ by setting $X_i=(x_{(i-1)\cdot s+1},\ldots,x_{i\cdot s})$. View each $i_j$ as a $\log n$-bit string with prefix $p_j\in\bits^{\log n-\log s}$ and a suffix $s_j\in\bits^{\log s}$. This is just relabeling wires, so it requires no circuitry.
\item Compute $(X_{p_1},\ldots,X_{p_q})\leftarrow C_n(X,(p_1,\ldots,p_q))$. This requires circuitry of size $S$ and depth $D$.
\item For each $j\in[q]$ in parallel, compute $x_{i_j}=\Sel_{\{0,1\}}^{s\to1}(X_{p_j},(s_j))$. By Proposition 2, for each $j\in[q]$ this requires circuitry of size $O(s)$ and depth $O(\log s)$, for a total circuitry size of $O(q\cdot s)$ and depth $O(\log s)$.
\item Output $(x_{i_1},\ldots,x_{i_q})$.
\end{enumerate}
\end{proof}
We can now prove our main theorem, which we recall here for convenience.
\begin{originalresult}{Main Theorem}
For all $n,q\in\ZZ^+$ there is a Boolean circuit computing $\Sel^{n\to q}$ of size $O(n+q\cdot\log^3(n))$ and depth $O(\log(n+q))$.
\end{originalresult}
\begin{proof}
Set $s(n)=\log^2n$ and let $\Sigma=\bits^s$. By Corollary 23, there exists a log-space uniform circuit for $\Sel_\Sigma^{n/s\to q}$ with size $O(s\cdot n/s+q\cdot s\cdot\log n)=O(n+q\log^3n)$ and depth $O(\log(n/s))=O(\log n)$. Applying Lemma 27 to this circuit yields a log-space uniform circuit for $\Sel_{\{0,1\}}^{n\to q}$ with size $O(n+q\cdot(\log^3n+s))=O(n+q\cdot\log^3n)$ and depth $O(\log n+\log s)=O(\log n)$.
\end{proof}
